\documentclass[11pt]{article}

\usepackage[margin=1in]{geometry}
\usepackage{amsmath, amssymb, mathtools}
\usepackage{booktabs}
\usepackage{array}
\usepackage{longtable}
\usepackage{siunitx}
\usepackage{enumitem}
\usepackage[hidelinks]{hyperref}

\sisetup{
  per-mode=symbol,
  group-separator={,}
}

\title{Stage 1 Mathematical Formulation for the Blown-Wing Concept Optimizer}
\author{AA146 Capstone Optimizer Notes}
\date{\today}

\begin{document}
\maketitle

\begin{abstract}
This document formalizes the current Stage 1 implementation of the blown-wing concept optimizer in the repository. The formulation corresponds to the present code path rather than to any stale or previously generated CSV artifact. Stage 1 is a coarse, architecture-level screening tool that evaluates distributed-propulsion configurations for a fixed-span wing using a surrogate propeller model, a two-zone blown/unblown wing model, a bisection-based operating-point solve, and a three-objective Pareto filter. The purpose of Stage 1 is not final flight-vehicle certification or high-fidelity aerodynamic closure; rather, it is to eliminate weak architectures and pass viable configurations to higher-fidelity downstream analysis.
\end{abstract}

\section{Problem Definition}

Stage 1 screens a discrete propulsion design space while holding the baseline wing span and baseline Stage 1 chord fixed. The screened design variables are
\begin{equation}
\mathbf{x} =
\left\{
N_p,\;
D,\;
\frac{P}{D},\;
f_{\mathrm{prop}}
\right\},
\end{equation}
where \(N_p\) is the propeller count, \(D\) is propeller diameter, \(P/D\) is propeller pitch ratio, and \(f_{\mathrm{prop}}\) is a coarse propeller-family label in the set
\(\{\texttt{high\_thrust}, \texttt{balanced}, \texttt{cruise}\}\).

Stage 1 then solves the low-speed and cruise RPMs internally. Thus RPM is an operating-state variable, not an outer-loop design variable.

\section{Fixed Inputs and Assumptions}

\subsection{Mission and mass assumptions}

\begin{table}[h]
\centering
\begin{tabular}{llc}
\toprule
Symbol & Quantity & Value \\
\midrule
\(m_{\mathrm{gross}}\) & Gross flight mass & \SI{5.0}{kg} \\
\(m_{\max}\) & Maximum built mass & \SI{5.0}{kg} \\
\(m_{\mathrm{fixed}}\) & Fixed non-propulsion mass & \SI{2.0}{kg} \\
\(V_{\mathrm{batt}}\) & Battery voltage & \SI{14.8}{V} \\
\(\epsilon_{\mathrm{batt}}\) & Battery specific energy & \SI{180}{Wh/kg} \\
\(f_{\mathrm{reserve}}\) & Battery reserve fraction & 0.20 \\
\(t_{\mathrm{climb}}\) & Low-speed energy segment & \SI{1.5}{min} \\
\(t_{\mathrm{loiter}}\) & Loiter duration & \SI{18.0}{min} \\
\(V_{\infty,\mathrm{low}}\) & Low-speed flight condition & \SI{4.0}{m/s} \\
\(V_{\infty,\mathrm{cruise}}\) & Cruise flight condition & \SI{10.0}{m/s} \\
\bottomrule
\end{tabular}
\caption{Global mission and mass assumptions used by Stage 1.}
\end{table}

\subsection{Atmospheric assumptions}

\begin{table}[h]
\centering
\begin{tabular}{llc}
\toprule
Symbol & Quantity & Value \\
\midrule
\(\rho\) & Air density & \SI{1.225}{kg/m^3} \\
\(\mu\) & Dynamic viscosity & \SI{1.81e-5}{Pa.s} \\
\(g\) & Gravity & \SI{9.80665}{m/s^2} \\
\(a\) & Speed of sound & \SI{343.0}{m/s} \\
\bottomrule
\end{tabular}
\caption{Atmospheric constants used throughout Stage 1.}
\end{table}

\subsection{Baseline wing and aerodynamic assumptions}

\begin{table}[h]
\centering
\begin{tabular}{llc}
\toprule
Symbol & Quantity & Value \\
\midrule
\(b\) & Span & \SI{2.0}{m} \\
\(c\) & Baseline Stage 1 chord & \SI{0.35}{m} \\
\(S\) & Baseline Stage 1 wing area & \SI{0.70}{m^2} \\
\(AR\) & Baseline aspect ratio & 5.714 \\
\(e\) & Oswald efficiency factor & 0.6591 \\
\(C_{L,\max,\mathrm{flapped}}\) & Baseline flapped section lift ceiling & 2.2 \\
\(C_{L,\mathrm{ceil}}\) & Enforced wing-system lift ceiling & 2.0 \\
\(C_{L,\max,\mathrm{clean}}\) & Clean section lift ceiling & 1.4 \\
\(C_{D0,\mathrm{flapped}}\) & Flapped profile drag parameter & 0.11 \\
\(C_{D0,\mathrm{clean}}\) & Clean profile drag parameter & 0.05 \\
\(f_{\mathrm{trim}}\) & Trim drag multiplier & 1.05 \\
\(f_{\mathrm{blown}}\) & Blown profile drag multiplier & 1.05 \\
\(f_{\mathrm{unblown}}\) & Unblown profile drag multiplier & 1.0 \\
\(k_{\mathrm{exp}}\) & Slipstream span expansion factor & 0.8 \\
\bottomrule
\end{tabular}
\caption{Baseline wing and aerodynamic constants.}
\end{table}

\subsection{Packing assumptions}

\begin{table}[h]
\centering
\begin{tabular}{llc}
\toprule
Symbol & Quantity & Value \\
\midrule
\(w_f\) & Fuselage width & \SI{0.20}{m} \\
\(c_f\) & Fuselage-to-inboard-prop edge clearance & \SI{1.0}{in} \\
\(c_p\) & Inter-prop edge clearance & \SI{1.0}{in} \\
\(c_t\) & Wingtip edge margin & \SI{1.5}{in} \\
\bottomrule
\end{tabular}
\caption{Stage 1 prop-packing assumptions.}
\end{table}

\subsection{Electrical, solver, and limit assumptions}

\begin{table}[h]
\centering
\begin{tabular}{llc}
\toprule
Symbol & Quantity & Value \\
\midrule
\(\eta_{\mathrm{chain}}\) & Electrical chain efficiency & 0.72 \\
\(P_{\mathrm{av}}\) & Avionics power & \SI{10}{W} \\
\(\gamma_{\mathrm{low}}\) & Low-speed thrust margin multiplier & 1.05 \\
\(\gamma_{\mathrm{cruise}}\) & Cruise thrust margin multiplier & 1.02 \\
\(M_{\mathrm{tip,max}}\) & Tip Mach limit & 0.55 \\
\(\mathrm{RPM}_{\mathrm{low}}\) bounds & Low-speed RPM bracket & \([4000,\,14000]\) \\
\(\mathrm{RPM}_{\mathrm{cruise}}\) bounds & Cruise RPM bracket & \([3000,\,11000]\) \\
\(\Delta \mathrm{RPM}_{\mathrm{tol}}\) & RPM bisection tolerance & \SI{5}{rpm} \\
\(\left(T/A\right)_{\min}\) & Minimum disk loading & \SI{25}{N/m^2} \\
\(\left(T/A\right)_{\max}\) & Maximum disk loading & \SI{500}{N/m^2} \\
\bottomrule
\end{tabular}
\caption{Electrical and solver limits enforced in Stage 1.}
\end{table}

\section{Stage 1 Design Space}

The present Stage 1 sweep is
\begin{align}
N_p &\in \{4,6,8,10,12,14,16\}, \\
D &\in \{4.0,4.5,5.0,5.5,6.0,6.5,7.0,7.5,8.0,9.0,10.0\}\;\mathrm{in}, \\
\frac{P}{D} &\in \{0.40,0.50,0.60,0.70,0.80,0.90\}, \\
f_{\mathrm{prop}} &\in \{\texttt{high\_thrust},\texttt{balanced},\texttt{cruise}\}.
\end{align}
Therefore the Stage 1 discrete candidate count is
\begin{equation}
7 \times 11 \times 6 \times 3 = 1386.
\end{equation}

\section{Propeller Surrogate Model}

\subsection{Family parameters}

\begin{table}[h]
\centering
\begin{tabular}{lccc}
\toprule
Family & \(C_{T,\mathrm{static,base}}\) & \(C_{P,\mathrm{static,base}}\) & \(J_{\mathrm{crit}}\) \\
\midrule
\texttt{high\_thrust} & 0.165 & 0.092 & 0.55 \\
\texttt{balanced} & 0.135 & 0.076 & 0.75 \\
\texttt{cruise} & 0.110 & 0.064 & 0.95 \\
\bottomrule
\end{tabular}
\caption{Stage 1 coarse prop-family constants.}
\end{table}

\subsection{Pitch-ratio scaling}

Let
\begin{equation}
n = \frac{\mathrm{RPM}}{60}
\end{equation}
be rotational speed in revolutions per second. The advance ratio is
\begin{equation}
J = \frac{V}{nD}.
\end{equation}
With reference pitch ratio \((P/D)_{\mathrm{ref}} = 0.65\), Stage 1 defines
\begin{align}
C_{T,\mathrm{static}} &= C_{T,\mathrm{static,base}}
\left(\frac{P/D}{0.65}\right)^{0.35}, \\
C_{P,\mathrm{static}} &= C_{P,\mathrm{static,base}}
\left(\frac{P/D}{0.65}\right)^{0.80}.
\end{align}

\subsection{Operating-point decay with advance ratio}

The thrust and power factors are
\begin{align}
f_T(J) &= \max\left(0,\; 1 - \left(\frac{J}{J_{\mathrm{crit}}}\right)^{1.45}\right), \\
f_P(J) &= \max\left(0.10,\; 1 - \left(\frac{J}{1.15\,J_{\mathrm{crit}}}\right)^{1.60}\right).
\end{align}
Thus,
\begin{align}
C_T &= C_{T,\mathrm{static}}\,f_T\,f_{Re}, \\
C_P &= C_{P,\mathrm{static}}\,(0.25 + 0.75 f_P)\,f_{Re},
\end{align}
where \(f_{Re}\) is the blade Reynolds derating defined below.

\subsection{Blade Reynolds derating}

At the \(0.75R\) station, Stage 1 computes
\begin{align}
r_{0.75} &= 0.75 \frac{D}{2}, \\
V_\theta &= 2\pi r_{0.75} n, \\
V_{\mathrm{res}} &= \sqrt{V_\theta^2 + V_\infty^2}, \\
c_{\mathrm{blade}} &= 0.10 D, \\
Re_{\mathrm{blade}} &= \frac{\rho V_{\mathrm{res}} c_{\mathrm{blade}}}{\mu}.
\end{align}
The Reynolds penalty is
\begin{equation}
f_{Re} = \max\left(0.60,\; \min\left[1,\;
\left(\frac{Re_{\mathrm{blade}}}{80000}\right)^{0.30}\right]\right).
\end{equation}

\subsection{Thrust, power, torque, and tip Mach}

Per propeller,
\begin{align}
T_{\mathrm{per}} &= C_T \rho n^2 D^4, \\
P_{\mathrm{shaft,per}} &= C_P \rho n^3 D^5.
\end{align}
For \(N_p\) identical propellers,
\begin{align}
T_{\mathrm{tot}} &= N_p T_{\mathrm{per}}, \\
P_{\mathrm{shaft,tot}} &= N_p P_{\mathrm{shaft,per}}, \\
P_{\mathrm{elec,tot}} &= \frac{P_{\mathrm{shaft,tot}}}{\eta_{\mathrm{chain}}} + P_{\mathrm{av}}.
\end{align}
Per-prop torque is
\begin{equation}
\tau_{\mathrm{per}} = \frac{P_{\mathrm{shaft,per}}}{2\pi n}.
\end{equation}
Tip speed and tip Mach are
\begin{align}
V_{\mathrm{tip}} &= \sqrt{(\pi D n)^2 + V_\infty^2}, \\
M_{\mathrm{tip}} &= \frac{V_{\mathrm{tip}}}{a}.
\end{align}

\section{Packing and Blown Span Fraction}

Only even propeller counts are allowed. For each half-wing, the first and last prop centers must satisfy the fuselage and tip constraints. If the required center-to-center packing distance exceeds the available half-span, the candidate is rejected.

For a given placed layout, the blown span fraction is formed by unioning the spanwise slipstream intervals. Each prop creates a blown span interval of width
\begin{equation}
w_{\mathrm{blown}} = k_{\mathrm{exp}} D,
\end{equation}
with \(k_{\mathrm{exp}} = 0.8\). Overlapping intervals are merged, and then
\begin{equation}
\eta_b = \frac{\text{merged blown span}}{b}.
\end{equation}

\section{Slipstream Velocity Model}

The total disk area is
\begin{equation}
A_{\mathrm{disk}} = N_p \pi \left(\frac{D}{2}\right)^2.
\end{equation}
Stage 1 uses an ideal actuator-disk relation:
\begin{align}
v_i &= \sqrt{\frac{T_{\mathrm{tot}}}{2 \rho A_{\mathrm{disk}}}}, \\
V_{\mathrm{eff}} &= V_\infty + 2 v_i.
\end{align}
The inverse form is used when the code needs the thrust required to generate a target blown velocity \(V_{\mathrm{eff,req}}\):
\begin{equation}
T_{\mathrm{req}}(V_{\mathrm{eff,req}})
=
2\rho A_{\mathrm{disk}}
\left(
\frac{V_{\mathrm{eff,req}} - V_\infty}{2}
\right)^2.
\end{equation}

\section{Low-Speed Two-Zone Blown-Wing Model}

\subsection{Area partition}

At low speed, the wing is partitioned into
\begin{align}
S_u &= (1-\eta_b)S, \\
S_b &= \eta_b S.
\end{align}
The corresponding dynamic pressures are
\begin{align}
q_\infty &= \frac{1}{2}\rho V_{\infty,\mathrm{low}}^2, \\
q_b &= \frac{1}{2}\rho V_{\mathrm{eff}}^2.
\end{align}

\subsection{Momentum-coefficient lift augmentation}

Stage 1 defines a wing-reference momentum coefficient
\begin{equation}
C_\mu = \frac{T_{\mathrm{tot}}}{q_\infty S}.
\end{equation}
The blown-section lift increment is
\begin{equation}
\Delta C_{L,\mu} =
\begin{cases}
0, & C_\mu < 0.02, \\
1.8 \sqrt{C_\mu}, & C_\mu \ge 0.02.
\end{cases}
\end{equation}
The section ceilings are then
\begin{align}
C_{L,\max,b} &= \min\left(C_{L,\max,\mathrm{flapped}} + \Delta C_{L,\mu},\; C_{L,\mathrm{ceil}}\right), \\
C_{L,\max,u} &= \min\left(C_{L,\max,\mathrm{flapped}},\; C_{L,\mathrm{ceil}}\right).
\end{align}
Since \(C_{L,\mathrm{ceil}}=2.0\), both blown and unblown sections are capped at \(2.0\) in the current implementation.

\subsection{Lift partition and required blown velocity}

The Stage 1 reference lift coefficient requirement is
\begin{equation}
C_{L,\mathrm{req,ref}} = \frac{W}{q_\infty S},
\qquad
W = m_{\mathrm{gross}} g.
\end{equation}
The unblown portion can support at most
\begin{equation}
L_{u,\max} = q_\infty S_u C_{L,\max,u}.
\end{equation}
The remaining required lift is
\begin{equation}
L_{\mathrm{rem}} = W - L_{u,\max}.
\end{equation}
If \(L_{\mathrm{rem}} \le 0\), then
\begin{equation}
V_{\mathrm{eff,req}} = V_{\infty,\mathrm{low}}.
\end{equation}
Otherwise,
\begin{align}
q_{b,\mathrm{req}} &= \frac{L_{\mathrm{rem}}}{S_b C_{L,\max,b}}, \\
V_{\mathrm{eff,req}} &= \sqrt{\frac{2 q_{b,\mathrm{req}}}{\rho}},
\end{align}
with \(V_{\mathrm{eff,req}} \ge V_{\infty,\mathrm{low}}\).

\subsection{Low-speed section lift requirements}

If \(L_{\mathrm{rem}} \le 0\), then the wing is low-speed feasible without requiring the blown section to carry load. Otherwise, the blown-section lift requirement is
\begin{equation}
C_{L,\mathrm{req},b} = \frac{L_{\mathrm{rem}}}{q_b S_b},
\end{equation}
and low-speed lift feasibility requires
\begin{equation}
C_{L,\mathrm{req},b} \le C_{L,\max,b}.
\end{equation}

\subsection{Low-speed drag}

The induced-drag factor is
\begin{equation}
k = \frac{1}{\pi e AR}.
\end{equation}
The low-speed section drag coefficients are
\begin{align}
C_{D,u} &= f_{\mathrm{trim}}
\left(
f_{\mathrm{unblown}} C_{D0,\mathrm{flapped}} + k C_{L,u}^2
\right), \\
C_{D,b} &= f_{\mathrm{trim}}
\left(
f_{\mathrm{blown}} C_{D0,\mathrm{flapped}} + k C_{L,b}^2
\right).
\end{align}
Thus the total low-speed drag is
\begin{equation}
D_{\mathrm{low}} = q_\infty S_u C_{D,u} + q_b S_b C_{D,b}.
\end{equation}

\section{Cruise Model}

At cruise, Stage 1 uses a clean-wing drag buildup. The cruise lift coefficient requirement is
\begin{equation}
C_{L,\mathrm{cruise}} = \frac{W}{q_{\mathrm{cruise}} S},
\qquad
q_{\mathrm{cruise}} = \frac{1}{2}\rho V_{\infty,\mathrm{cruise}}^2.
\end{equation}
The cruise drag coefficient is
\begin{equation}
C_{D,\mathrm{cruise}} = C_{D0,\mathrm{clean}} + k C_{L,\mathrm{cruise}}^2,
\end{equation}
and cruise drag is
\begin{equation}
D_{\mathrm{cruise}} = q_{\mathrm{cruise}} S C_{D,\mathrm{cruise}}.
\end{equation}
Cruise lift feasibility requires
\begin{equation}
C_{L,\mathrm{cruise}} \le C_{L,\max,\mathrm{clean}}.
\end{equation}

\section{Operating-Point Solve}

\subsection{Low-speed RPM solve}

Stage 1 solves the minimum low-speed RPM that satisfies both the drag requirement and the blown-velocity requirement. The low-speed thrust requirement is
\begin{equation}
T_{\mathrm{req,low}} =
\max\left(
\gamma_{\mathrm{low}} D_{\mathrm{low}},\;
T_{\mathrm{req}}(V_{\mathrm{eff,req}})
\right).
\end{equation}
The low-speed residual is
\begin{equation}
R_{\mathrm{low}}(\mathrm{RPM}) =
T_{\mathrm{tot}}(\mathrm{RPM}) - T_{\mathrm{req,low}}.
\end{equation}
Stage 1 uses a bisection solve over the bracket \([4000,14000]\) rpm.

\subsection{Cruise RPM solve}

The cruise thrust requirement is
\begin{equation}
T_{\mathrm{req,cruise}} = \gamma_{\mathrm{cruise}} D_{\mathrm{cruise}},
\end{equation}
and the cruise residual is
\begin{equation}
R_{\mathrm{cruise}}(\mathrm{RPM}) =
T_{\mathrm{tot}}(\mathrm{RPM}) - T_{\mathrm{req,cruise}}.
\end{equation}
This residual is solved over \([3000,11000]\) rpm.

\section{Motor and Mass Inference}

\subsection{Motor requirements}

Stage 1 does not select a catalog motor directly. It infers motor requirements from the solved operating points:
\begin{align}
I_{\mathrm{motor}} &=
\frac{\max(P_{\mathrm{elec,low}}, P_{\mathrm{elec,cruise}})}
{V_{\mathrm{batt}} N_p}, \\
k_V &\approx \frac{\mathrm{RPM}_{\mathrm{low}}}{0.85 V_{\mathrm{batt}}}.
\end{align}

\subsection{Propulsion mass fits}

The Stage 1 propulsion-mass closure uses fitted power laws from local CSV databases:
\begin{align}
m_{\mathrm{prop}}[\mathrm{g}] &=
0.10287476896145267 \;
D_{\mathrm{in}}^{2.503890170879269}
\left(1 + 0.014466435874827321 \frac{P}{D}\right), \\
m_{\mathrm{motor}}[\mathrm{g}] &=
0.39551092994895437 \;
P_{\mathrm{shaft,peak}}^{0.8602954664135029}, \\
m_{\mathrm{ESC}}[\mathrm{g}] &=
0.3553247179813882 \;
I_{\mathrm{motor}}^{1.1780662036609257}.
\end{align}

\subsection{Battery mass}

The mission energy model is
\begin{equation}
E_{\mathrm{usable}} =
P_{\mathrm{low}} \frac{t_{\mathrm{climb}}}{60}
+
P_{\mathrm{cruise}} \frac{t_{\mathrm{loiter}}}{60}.
\end{equation}
The required pack energy and battery mass are
\begin{align}
E_{\mathrm{pack}} &= E_{\mathrm{usable}}(1 + f_{\mathrm{reserve}}), \\
m_{\mathrm{battery}} &= \frac{E_{\mathrm{pack}}}{\epsilon_{\mathrm{batt}}}.
\end{align}

The total built mass is
\begin{equation}
m_{\mathrm{built}} = m_{\mathrm{fixed}} + m_{\mathrm{propulsion}} + m_{\mathrm{battery}},
\end{equation}
and Stage 1 requires
\begin{equation}
m_{\mathrm{built}} \le m_{\max}.
\end{equation}

\section{Feasibility Conditions}

A Stage 1 candidate is feasible only if all of the following conditions are satisfied:
\begin{align}
\text{packing margin} &\ge 0, \\
\text{low-speed lift feasible} &\equiv C_{L,\mathrm{req},b} \le C_{L,\max,b}, \\
\text{cruise lift feasible} &\equiv C_{L,\mathrm{cruise}} \le C_{L,\max,\mathrm{clean}}, \\
V_{\mathrm{eff}} - V_{\mathrm{eff,req}} &\ge 0, \\
T_{\mathrm{tot,low}} - T_{\mathrm{req,low}} &\ge 0, \\
T_{\mathrm{tot,cruise}} - T_{\mathrm{req,cruise}} &\ge 0, \\
M_{\mathrm{tip,low}} &\le M_{\mathrm{tip,max}}, \\
M_{\mathrm{tip,cruise}} &\le M_{\mathrm{tip,max}}, \\
\mathrm{RPM}_{\mathrm{low}} - \mathrm{RPM}_{\mathrm{cruise}} &\ge 0, \\
\left(T/A\right)_{\min} \le \frac{T_{\mathrm{tot,low}}}{A_{\mathrm{disk}}} &\le \left(T/A\right)_{\max}, \\
m_{\mathrm{built}} &\le m_{\max}.
\end{align}

\section{Stage 1 Objective Set}

Stage 1 does not collapse the design space to a single scalar objective. Instead, it computes a three-objective Pareto front:
\begin{align}
\min \; P_{\mathrm{low}}, \qquad
\min \; E_{\mathrm{loiter}}, \qquad
\min \; m_{\mathrm{propulsion}}.
\end{align}
The resulting nondominated set is then passed to later stages.

\section{Outputs Passed Downstream}

For every feasible candidate, Stage 1 records:
\begin{itemize}[leftmargin=1.5em]
\item solved low-speed and cruise RPM,
\item low-speed and cruise thrust and drag,
\item low-speed actual and required blown velocity,
\item blown span fraction,
\item low-speed and cruise \(C_L\) requirements,
\item disk loading and momentum coefficient,
\item per-prop torque and peak shaft power,
\item inferred motor current and \(k_V\),
\item propulsion mass, battery mass, and built-mass margin,
\item all constraint margins.
\end{itemize}

\section{Modeling Limits}

The present Stage 1 model has several intentional simplifications:
\begin{itemize}[leftmargin=1.5em]
\item propeller aerodynamics are surrogate-based, not yet driven by full measured UIUC \(C_T\)--\(C_P\) maps;
\item the slipstream model is actuator-disk based and span-averaged rather than fully resolved;
\item the low-speed wing is treated as a two-zone blown/unblown system rather than a continuously varying spanwise flowfield;
\item Stage 1 wing geometry is fixed;
\item structural sizing, control-surface sizing, and coupled stability analysis are deferred to later stages.
\end{itemize}

\section{Interpretation}

Stage 1 should therefore be read as a concept downselection model. It answers the question:
\emph{which distributed-propulsion architectures can plausibly satisfy the low-speed blown-lift requirement, the cruise requirement, the propulsion-mass closure, and the packing limits under a common set of conservative screening assumptions?}

This is precisely the role a front-end architecture optimizer should play before higher-fidelity geometry refinement and more strongly coupled aerodynamic modeling are introduced.

\end{document}
